\chapter{核心思想与设计原则}

\section{本章作用}

前四章已经分别说明了 Pebble 的问题域、项目主张、背景冲突与术语边界。本章进一步回答一个更基础的问题：为什么 Pebble 必须围绕关系上下文组织为一条连续支持链路，以及哪些原则在后续方法、流程与系统设计中不应被破坏。

这一章的任务是给出一套足以约束后续设计判断的核心思想。缺少这一层时，后文关于解码、陪练、稳态恢复、关系档案与洞察累积的设计容易被误读为若干分散功能的叠加；核心思想被固定后，后续章节中的模块划分、数据组织和交互路径就会呈现为围绕同一目标展开的结构安排。

\section{核心思想：以关系上下文组织连续支持}

Pebble 的核心思想可以概括为一句话：系统围绕具体关系对象维护可持续更新的上下文，并把这一上下文注入用户每一次理解、练习、恢复与复盘过程，使支持能力随着互动历史逐步变得更贴近真实处境。

这一思想包含四个相互依赖的判断。

\begin{itemize}
  \item 对复杂对话的解释不能脱离关系背景。同一句表达在亲密关系、家庭关系或职场关系中往往具有完全不同的压力结构、风险信号与回应空间。
  \item 用户需求体现为连续支持。用户需要先理解话语，再决定如何回应，在情绪负荷过高时恢复状态，并在后续互动中持续校正判断。
  \item 关系洞察应以渐进累积方式形成。系统对关系对象的理解来自多轮互动、场景补充与重复出现的行为线索，而不是一次输入后的永久结论。
  \item 产品价值来自闭环，而不是单点准确率。即使某一次分析文本的结果看似合理，若无法被带入陪练、记录和后续适配，系统仍难以形成稳定的长期价值。
\end{itemize}

因此，Pebble 的“核心思想”并非某一种单独的模型能力，也不是一条孤立的界面流程。它是一种组织原则：把关系对象作为解释中心，把用户支持链路作为设计中心，把长期累积作为系统演化中心。

\section{闭环逻辑：从单次解释走向持续洞察}

在这一思想下，Pebble 的能力结构可被理解为一条由四段组成的闭环。第一段负责理解当前输入，第二段负责把理解结果转化为可练习的应对方案，第三段负责在用户状态失稳时提供最小可用的恢复支撑，第四段负责把新一轮互动中的观察沉淀回关系档案，为后续解释提供更稳定的上下文条件。

这一闭环的关键不在于模块数量，而在于信息如何回流。系统若只停留在“输入一句话，返回一段解释”的层面，就只能生成局部答案；系统若能把关系背景、观察记录和用户选择持续带回下一轮分析，才有条件形成逐步增强的情境适配能力。由此，Pebble 的长期价值不应只用单次回答的表面可读性来衡量，还应考察它是否使用户在多轮互动中逐渐建立更稳的判断框架和边界策略。

从报告结构看，这一闭环同时解释了后续章节的组织方式。

\begin{itemize}
  \item 方法与建模章节会说明系统如何把关系对象、对话输入与观察记录组织成可被持续调用的分析条件。
  \item 流程与执行路径章节会说明用户如何在页面之间完成理解、练习、恢复与回流。
  \item 系统与数据设计章节会说明哪些数据必须长期维护，哪些状态只在局部交互中存在。
  \item 示例、结果与验证章节会说明系统价值应如何通过闭环效果而非单点表现来理解。
\end{itemize}

\section{设计原则}

核心思想若要在实现中保持稳定，至少需要六条设计原则作为约束。它们并非文案层面的偏好，而是后续架构与交互设计的判断准则。

表~\ref{tab:pebble-core-principles} 汇总了上述原则及其约束含义。

\begin{table}[H]
  \centering
  \caption{Pebble 核心设计原则及其在系统中的约束含义。}
  \label{tab:pebble-core-principles}
  \begin{tabularx}{\textwidth}{>{\raggedright\arraybackslash}p{2.8cm} >{\raggedright\arraybackslash}X >{\raggedright\arraybackslash}X}
    \toprule
    原则 & 含义 & 对后续设计的约束 \\
    \midrule
    关系上下文优先 & 解释与建议应围绕具体关系对象组织，而非只依据输入文本表面语义 & 后续分析服务、状态管理与页面流程都应允许关系信息进入主路径 \\
    先理解后应对 & 系统先帮助用户辨识语气、潜台词与风险，再进入回应建议与练习阶段 & 解码输出应承担中介角色，不能跳过理解层直接给模板回复 \\
    渐进累积 & 关系洞察来自多轮互动中的重复信号与上下文补充 & 数据层应支持持续更新和回看，避免把一次分析写成永久性定论 \\
    去标签化 & 系统用观察性、描述性语言表述模式，不对关系对象做武断定性 & 洞察表达应强调行为线索与证据，不把诊断性词汇当成默认输出 \\
    用户主权 & 用户对关系档案、洞察记录与策略采纳保持最终决定权 & 交互设计应保留查看、编辑、忽略或删除的空间，避免系统替用户下判断 \\
    边界克制与隐私优先 & 系统只收集完成支持所需的最小信息，并对敏感关系信息保持克制处理 & 后续设计应避免不必要的数据扩张，并明确区分即时分析与长期沉淀的边界 \\
    \bottomrule
  \end{tabularx}
\end{table}

上述原则可以进一步压缩为三组更高层的约束关系。

\begin{itemize}
  \item 以关系上下文优先、先理解后应对为代表的原则，约束系统应如何组织分析入口。
  \item 以渐进累积、去标签化为代表的原则，约束系统应如何形成和表达关系洞察。
  \item 以用户主权、边界克制与隐私优先为代表的原则，约束系统在提供支持时不得越过的责任边界。
\end{itemize}

\section{原则如何落到系统面}

这些原则的意义在于，它们能够直接解释若干后续设计选择为何成立。

\begin{itemize}
  \item 若坚持关系上下文优先，那么“当前关系对象”的选择与调用就不能只是页面附属信息，而应成为解码与建议路径中的主条件。
  \item 若坚持先理解后应对，那么系统输出就应先呈现表层语义、潜台词与情绪风险，再进入回复建议与练习动作。
  \item 若坚持渐进累积与去标签化，那么关系档案中沉淀的内容就应以观察记录、证据线索与可更新摘要为中心，而非一次性贴上结论性标签。
  \item 若坚持用户主权与边界克制，那么 Pebble 的角色就应被限定为支持性工具，系统可以提供洞察、建议与训练环境，但不能越过用户的价值判断与现实决策责任。
\end{itemize}

由此可见，第五章提出的是一组能够约束交互流程、服务编排、数据模型与内容表达的基础规则。后续章节若出现任何局部设计选择，读者都可以回到本章检验一个问题：该选择是否仍然服务于“关系上下文驱动的连续支持”这一核心思想，并且没有破坏本章确立的设计原则。

\section{本章小结}

本节读完后，读者应能解释为什么 Pebble 必须围绕关系上下文组织为连续支持链路，并列举至少三条约束后续设计选择的核心原则。

\reportadd{
\section{关键设计决策的权衡}

Pebble 的若干核心设计并非唯一可行路径。本节明确三项关键决策的主要权衡，以及替代方案被排除的主要理由。

表~\ref{tab:pebble-tradeoffs} 对上述三项决策做了进一步压缩。

\begin{table}[H]
  \centering
  \caption{Pebble 核心设计决策的权衡对照。}
  \label{tab:pebble-tradeoffs}
  \begin{tabularx}{\textwidth}{>{\raggedright\arraybackslash}p{2.2cm} >{\raggedright\arraybackslash}p{2.8cm} >{\raggedright\arraybackslash}X >{\raggedright\arraybackslash}X}
    \toprule
    决策 & 选定方案 & 主要收益 & 主要代价与被排除的替代方案 \\
    \midrule
    分析入口的组织方式 & 关系上下文优先：以关系对象为中心组织分析 & 同一句话在不同关系中获得差异化解释；长期洞察可渐进累积 & 需维护额外数据层；被排除的替代是纯文本单轮分析——更简单即时，但失去情境适配与长期记忆能力，与产品核心价值主张冲突 \\
    系统架构深度 & Next.js BFF + PostgreSQL + Drizzle 六层全栈结构 & 支持数据沉淀、匿名/登录双轨、练习本持久化；确定性 fallback 可在服务端执行 & 引入数据库运维、迁移管理复杂度；被排除的替代是纯前端直接调用 LLM API——上线更快，但失去用户数据持久化、跨设备同步与精细权限控制 \\
    用户身份策略 & Cookie-based guest session + 可选 Supabase Auth 登录 & 降低首次体验门槛；匿名用户可先体验核心功能 & 数据模型需同时支持 guestSessionId 和 userId；被排除的替代是强制登录——数据模型更简单，但提高使用门槛，与"低摩擦支持"的产品目标冲突 \\
    \bottomrule
  \end{tabularx}
\end{table}
}
